% titel: Blatt 3, A3: Balken im Gleichgewicht
% skript: Kap. Statik, Drehmoment

% Drehmomenten-Gleichgewicht eines drehbar gelagerten Balkens.

Ein starrer Balken ist im Gleichgewicht, wenn die Summe aller Kräfte und
die Summe aller Drehmomente (Momente) verschwinden.

\begin{werkzeugkasten}[Werkzeuge]
\textbf{Drehmoment} einer Kraft $F$ um einen Drehpunkt:
\begin{equation}
  M=F_\perp\,r
\end{equation}
mit $F_\perp$ der zur radialen Richtung senkrechten Kraftkomponente und
$r$ dem Hebelarm (Abstand Drehpunkt--Angriffspunkt). Für eine
vertikal hängende Masse ist der Hebelarm die horizontale Distanz vom
Drehpunkt.

\textbf{Statisches Gleichgewicht} eines gelagerten Körpers:
\begin{equation}
  \sum \vec F=\vec 0,\qquad \sum M=0\quad(\text{um beliebigen Punkt}).
\end{equation}
Waehlt man den Drehpunkt als Lager, fällt die unbekannte Lagerkraft aus
der Momentengleichung.

\textbf{Winkel am geneigten Balken:} Ein Arm, der um den Winkel $\alpha$
gegen die Horizontale geneigt ist und die Länge $\ell$ hat, erzeugt
einen horizontalen Hebelarm $\ell\cos\alpha$; ein senkrecht abgeknicktes
Stuck der Länge $b$ gibt zusätzlich $b\sin\alpha$ (oder zieht den Arm
zurück, je nach Richtung).
\end{werkzeugkasten}

\begin{werkzeugkasten}[Vorgehen]
\begin{itemize}
  \item Skizze mit allen Kräfte (den \emph{senkrecht} hängenden
        Gewichten $m_1g$, $m_2g$ und der Lagerkraft im Drehpunkt).
  \item Momentengleichung um das Lager aufstellen; die Lagerkraft fällt
        heraus.
  \item Die Hebelarme als $\ell\cos\alpha \pm b\sin\alpha$ ausdrücken und
        nach $\alpha$ Auflösen.
\end{itemize}
\end{werkzeugkasten}

\begin{achtung}[Typische Fallen]
\begin{itemize}
  \item Vorzeichen der Momente (Drehsinn) konsequent festlegen; gegen den
        Uhrzeigersinn ist z. B. positiv.
  \item Für eine vertikal wirkende Kraft zählt nur die \emph{horizontale}
        Distanz vom Lager.
  \item Das abgeknickte Stück traegt zum Hebelarm bei; bei \emph{grossen}
        Winkeln kann $\cos\alpha$ stark schrumpfen.
\end{itemize}
\end{achtung}
